\section{In-Kernel Neural Policy Inference}
\label{sec:implementation}
\label{sec:in_kernel_inference}

RL-SchedExt는 Android 15의 Linux 6.1 Kernel에
Backport된 sched\_ext Framework를 기반으로 구현하였다.
학습된 Neural Student의 Runtime 실행 과정은
Runtime State Collection, Student Neural Policy Inference 및
eBPF Scheduling Decision으로 구성된다.

RL-SchedExt는 기존 sched\_ext의 Task Lifecycle 및 Dispatch 경로에
Native Kernel 기반의 State Collection과 Neural Network Inference를
통합하고, Task Selection을 위한 eBPF Callback을 추가하였다.
Scheduling Decision에 필요한 Feature Construction 및
Neural Network 연산은 Native Kernel에서 수행하며,
eBPF Program에는 각 Candidate의 Vruntime Lag와
Learned Vruntime Adjustment를 전달한다.
이를 통해 eBPF Verifier의 Program Complexity 및 Stack 제약에
대한 부담을 줄이면서 Scheduling Policy의 최종 Task Selection은
eBPF Program에서 수행할 수 있다.
별도의 User Space Inference Daemon과 Runtime 학습은 필요하지 않다.

본 구현의 Neural Policy는 CPU7의 Local Dispatch Queue (DSQ)를
대상으로 동작한다.
따라서 본 연구의 적용 범위는 하나의 CPU에서 경쟁하는
Runnable Task의 실행 순서를 결정하는 것이며,
CPU Selection 및 Load Balancing을 학습하는 것은 아니다.

% Implementation reference:
% android-gs-caimito-6.1-android15-qpr2-beta
% integration/ctx5-tiny-linear-v6-20261122
% aosp/kernel/sched/ext.c, scx_rl_ctx5_native.h,
% scx_rl_ctx5_native_history.h, scx_rl_ctx5_tiny.h
% aosp/tools/sched_ext/include/scx/rl_ctx5_export_core.h,
% rl_ctx5_tiny_infer.h, rl_ctx5_tiny.bpf.h
% aosp/tools/sched_ext/scx_simple.bpf.c

\subsection{Runtime State Collection}

\paragraph{Event-driven Runtime History.}
System Design에서 정의한 State Feature는 단일 Scheduling Event에서
모두 관측할 수 있는 정보가 아니다.
예를 들어 Wakeup Rate와 CPU Utilization은 일정 기간의 실행 이력이
필요하며, Current Wait는 Task가 Runnable 상태로 전환된 시점과
실제로 실행되는 시점 사이의 시간 정보를 필요로 한다.
따라서 RL-SchedExt는 Native Kernel에서 Task별 Runtime History를
관리하고, Scheduling Event가 발생할 때 관련 정보를 갱신한다.

Table~\ref{tab:runtime_events}는 State Collection을 위해 추가한
Native Kernel Event와 주요 갱신 정보를 나타낸다.
Task History는 PID와 Task의 생성 시각을 함께 사용하여 식별하며,
Task의 Scheduling Class 또는 CPU 소속이 변경되는 경우
기존 History의 유효성을 관리한다.
Native Policy Mode에서는 eBPF의 기존 History 갱신 경로를
사용하지 않아 동일 Event가 중복으로 반영되지 않도록 한다.

\begin{table}[htbp]
    \centering
    \small
    \renewcommand{\arraystretch}{1.15}
    \begin{tabularx}{\linewidth}{|p{0.24\linewidth}|X|}
        \hline
        \textbf{Scheduling Event} & \textbf{Runtime History Update} \\ \hline
        Runnable & Runnable transition, waiting start, wakeup history \\ \hline
        Running & Waiting time, execution start, context switch history \\ \hline
        Stopping & Execution accounting, preemption history, waiting state \\ \hline
        Quiescent & Sleep/block history, runnable state removal \\ \hline
        Weight Change & Scheduling weight and service accounting \\ \hline
        Reset & Task history initialization or invalidation \\ \hline
    \end{tabularx}
    \caption{Runtime history updates on scheduling events.}
    \label{tab:runtime_events}
\end{table}

CPU Utilization과 Runnable Ratio는 시간에 따른
Exponentially Weighted Moving Average (EWMA)를 이용하여 관리한다.
Wakeup 등의 빈도 정보는 제한된 크기의 Timestamp Ring을
이용하여 계산한다.
CPU Execution Time 및 Scheduling Weight를 기반으로
각 Task에 기대되는 CPU Service와 실제 제공된 CPU Service의
차이를 누적하여 Service Debt를 구성한다.
이와 같이 Event-driven Accounting을 사용함으로써
Scheduling Decision마다 과거의 전체 실행 이력을
다시 탐색할 필요가 없다.

\paragraph{Shadow Vruntime Accounting.}
sched\_ext에서 관리하는 Dispatch Queue의 Virtual Time은
CFS Scheduling Entity의 Vruntime과 동일한 Accounting 값이 아니다.
또한 Task가 sched\_ext Scheduling Class에서 실행되는 동안
CFS의 Vruntime이 기존 CFS 실행 경로와 동일하게
갱신된다고 가정할 수 없다.
따라서 RL-SchedExt는 CFS 기반 Scheduling Policy를 학습한
Neural Student의 입력과 Task Selection에 사용하기 위해
별도의 Shadow Vruntime을 관리한다.

Task $i$가 $\Delta t_i$ 동안 실행된 경우
Shadow Vruntime은 다음과 같이 갱신한다.

\begin{equation}
    V_i^{\mathrm{shadow}} \leftarrow
    V_i^{\mathrm{shadow}} +
    \Delta t_i \frac{w_{\mathrm{ref}}}{w_i}.
    \label{eq:shadow_vruntime}
\end{equation}

여기서 $w_i$는 sched\_ext Scheduling Weight이며,
$w_{\mathrm{ref}}$는 기준 Weight이다.
실행 시간에 따른 Weighted Accounting과 Wakeup 시의
Vruntime Placement를 별도로 관리하여
CFS의 주요 Virtual Runtime Dynamics를 모사한다.
Learned Vruntime Adjustment는 이러한 Accounting 과정에
반영하지 않는다.

Scheduling Decision 시점에는 Candidate Task Set의
최소 Shadow Vruntime을 기준으로 Vruntime Lag를 계산한다.

\begin{equation}
    \ell_i = V_i^{\mathrm{shadow}} -
    \min_{j\in\mathcal{C}} V_j^{\mathrm{shadow}}.
    \label{eq:runtime_vruntime_lag}
\end{equation}

이를 통해 절대적인 Vruntime 값 대신 Candidate Task 간의
상대적인 CPU Service 차이를 State Feature로 사용한다.
이 Shadow Accounting은 CFS의 주요 동작을 모사하는 것이며,
CFS의 전체 Scheduling Hierarchy와 동작을
완전히 동일하게 재현하는 것은 아니다.

\paragraph{Candidate Collection and State Construction.}
Task Selection 시점에는 Runqueue Lock을 유지한 상태에서
Local DSQ에 존재하는 Task 중 현재 CPU에서 실행 가능한
Candidate를 수집한다.
Scheduling Class, Runnable 상태 및 CPU Affinity를 확인하고,
Task Identity를 포함하는 Scheduling Context를 구성한다.
History가 유효하지 않거나 Candidate 수가 최대 16개를
초과하는 경우에는 일부 Candidate만을 이용하여
Neural Policy를 실행하지 않는다.

수집된 Candidate의 Runtime History를 현재 시각까지 갱신한 후,
13차원 Task Feature와 6차원 Global Feature를 추출한다.
Task Set의 Feature 합계, 최댓값 및 최소 Vruntime을 가지는
첫 번째와 두 번째 Candidate 등의 공통 정보는
Scheduling Decision마다 한 번만 계산한다.
이 정보를 각 Task의 State Construction에서 재사용하며,
Task별 상대 순위와 다른 Candidate와의 Feature 차이를
계산하여 38차원 Derived Context Feature를 구성한다.

최종 State는 Simulator에서 정의한 Android Feature Representation과
동일한 Feature 순서 및 Scaling을 사용하여
57차원의 Signed Q10 Vector로 변환한다.
State의 Feature Scaling과 Saturation은
Runtime에서 학습 시 사용한 입력 범위를 유지하기 위한 것이다.
이 과정에서 Application-level Deadline이나
Request Completion 정보는 사용하지 않는다.

\subsection{Student Neural Policy Inference}

\paragraph{Integer-only Neural Inference.}
Kernel Scheduling Context에서는 Floating-point 연산의 사용에
제약이 있으므로, 학습된 Neural Student를
Integer 연산 기반의 Network로 변환하였다.
학습 완료 후 Fully Connected Layer의 Weight를
Output Channel별 Symmetric INT8 형식으로 양자화하며,
Scale은 Q24, Bias와 Intermediate Activation은 Q10 형식으로
표현한다.
양자화된 Model Parameter와 Activation 계산에 필요한
정적 배열은 C Header 형태로 Kernel에 포함한다.

각 Fully Connected Layer의 Output Channel $o$는
다음과 같은 Fixed-point 연산으로 계산된다.

\begin{equation}
    y_o^{\mathrm{Q10}} =
    \operatorname{round}
    \left(
        \frac{
            m_o^{\mathrm{Q24}}\sum_j
            x_j^{\mathrm{Q10}}q_{o,j}
        }{2^{24}}
    \right)+b_o^{\mathrm{Q10}},
    \label{eq:student_integer_affine}
\end{equation}

여기서 $q_{o,j}$는 INT8 Weight,
$m_o^{\mathrm{Q24}}$는 Output Channel별 Scale이며,
$b_o^{\mathrm{Q10}}$는 Bias이다.
곱셈 및 누산에는 Signed 64-bit Integer를 사용하고,
Fixed-point 변환에는 Round-half-away-from-zero 규칙을 적용한다.
연산 중 Overflow가 발생한 경우에는 해당 결과를 사용하지 않는다.

Tanh Activation은 입력을 $[-4,4]$ 범위로 제한한 후,
Q10 간격으로 미리 계산된 8,193개의 값을 조회하여 수행한다.
따라서 Scheduling Hot Path에서 Floating-point 연산이나
Exponential Function 계산이 필요하지 않다.
이 Table은 Scheduling Policy가 아니라 Activation Function을
구현하기 위한 것으로, Neural Network의 Weight와
Set-context 연산은 그대로 유지된다.

\paragraph{Bounded Inference and Fast Path.}
DeepSets 구조에서는 각 Task의 Embedding을 계산한 후
전체 Candidate Task Set의 Context를 구성해야 한다.
이에 따라 Native Kernel의 Forward Pass를
Task Encoding, Mean/Max Pooling 및 Priority Head의
세 단계로 분리하였다.
전체 Candidate의 Embedding을 먼저 계산하여 저장하고,
Set Context는 한 번만 계산하여 각 Task의 Priority Head에서
공통으로 사용한다.

최대 Candidate 수를 16개로 제한하고,
Input State와 Intermediate Activation을 저장하기 위한
Scratch Buffer를 사전에 확보한다.
이를 통해 Runqueue Lock을 유지한 Scheduling Hot Path에서
동적 Memory Allocation을 수행하지 않는다.
Fully Connected Layer의 연산량은 Candidate 수 $N$에 대해
$1{,}720N$ MAC으로 증가하며,
Feature Construction의 Task별 상대 순위 계산에는
별도의 Candidate 간 비교가 필요하다.

또한 일반적인 Next-task Selection에서 Candidate가 하나뿐인 경우에는
Priority Adjustment가 Task Selection 결과를 변경할 수 없다.
따라서 단일 Candidate에 대해서는 State Construction 및
유효성 확인 이후 Neural Network의 Forward Pass를 생략하고,
0의 Adjustment를 전달하는 Fast Path를 구현하였다.
이는 Wakeup 평가나 Audit Mode까지 생략하는 방식은 아니다.

Priority Head의 출력은 $\pm 2$\,ms 범위로 제한하며,
Nanosecond 단위의 정수로 변환한다.
계산된 Adjustment와 Candidate별 Valid Flag를
Scheduling Context에 기록한 후 eBPF Program으로 전달한다.

\subsection{eBPF Scheduling Decision}

\paragraph{Task Selection Callback.}
기존 sched\_ext의 Local DSQ 기반 Task Selection 경로에
RL-SchedExt의 Task Selection Callback인
\texttt{rl\_choose}를 추가하였다.
Native Kernel은 다음 Task를 선택해야 하는 시점에
Candidate Collection과 Neural Inference를 수행하고,
\texttt{scx\_rl\_pick\_ctx}를 통해
Candidate의 Vruntime Lag, Learned Adjustment 및
유효성 정보를 eBPF Program에 전달한다.

eBPF Program은 전달된 Candidate에 대해
Equation~\ref{eq:urgency_aware_vruntime}에 대응하는
Scheduling Key를 계산한다.

\begin{equation}
    K_i = \ell_i + \Delta V_i^S.
    \label{eq:runtime_scheduling_key}
\end{equation}

여기서 $\ell_i$는 Equation~\ref{eq:runtime_vruntime_lag}에서
정의한 Vruntime Lag이다.
동일한 기준 Vruntime을 모든 Candidate에서 차감하므로
절대 Shadow Vruntime에 Adjustment를 적용하는 경우와
동일한 순서를 얻을 수 있다.
% 실제 코드의 tiny_corr_ns는 양수일 때 우선순위를 높이는
% correction이며, 본문의 Delta V와 부호가 반대이다.
% K_i = ctx5_native_lag_ns - tiny_corr_ns.

eBPF Callback은 ABI Version, Model Hash 및 Ready Flag를
검증한 후 Scheduling Key가 가장 작은 Candidate의 Index를
반환한다.
동일한 Key에 대해서는 기존 Candidate 순서를 유지하며,
Runtime Task Selection에는 Softmax와 Sampling이 필요하지 않다.
최종적으로 Native Kernel은 반환된 Index에 해당하는
Task의 실행 가능 여부를 검증하고 Dispatch를 수행한다.

\paragraph{Preserving CFS Scheduling Mechanisms.}
Neural Student는 어떤 Task를 다음에 실행할 것인지를 결정하지만,
현재 Task를 언제 선점할 것인지 또는 선택된 Task에
얼마나 긴 Time Slice를 부여할 것인지는 직접 결정하지 않는다.
이러한 요소까지 동시에 변경하면
Neural Policy에 의한 Task Ordering의 효과를 분리하기 어렵다.

따라서 RL-SchedExt는 sched\_ext 환경에 CFS 방식의
Time Slice 계산과 Preemption Mechanism을 구현하였다.
Time Slice는 Runnable Task 수와 Scheduling Weight의 합,
Scheduling Latency 및 Minimum Granularity를 기반으로 계산한다.
Task의 실행 시간이 계산된 Time Slice에 도달하는 시점에는
기존 High-resolution Scheduler Tick (HRTICK)을 활용하여
Rescheduling을 요청할 수 있도록 하였다.
또한 Runnable Task Set이 변경될 때 현재 Task의
Time Slice를 재평가한다.

Wakeup Preemption은 CFS 방식의 Vruntime 비교와
Wakeup Granularity를 이용하며,
Periodic Tick에서는 실제 실행 시간과 Ideal Time Slice,
다른 Runnable Task와의 Virtual Runtime 차이를 고려한다.
이러한 Preemption Mechanism이 Rescheduling을 요청하면
새로운 Scheduling Decision 시점에서 Neural Student가
Candidate Task의 실행 우선순위를 다시 계산한다.
즉, 평가에 사용한 CFS-style Wakeup 경로에서는
Wakeup 자체가 Neural Network를 통해 선점 여부를 결정하지 않는다.
Learned Adjustment는 Task Ordering에만 적용되며,
Weighted Vruntime Accounting이나 Time Slice의 크기를
직접 변경하지 않는다.

\paragraph{Candidate Validation and Fallback.}
Neural Policy를 Kernel Scheduling에 적용하기 위해서는
Inference 결과의 정확성뿐만 아니라
실행 가능한 Task가 선택되는지에 대한 검증이 필요하다.
RL-SchedExt는 Native State Snapshot과 eBPF Callback의
반환 결과를 검증하며,
Candidate 수 초과, 유효하지 않은 History,
Model ABI 또는 Hash 불일치, Integer Overflow 등의 경우에는
Neural Policy의 결정을 적용하지 않는다.
유효한 Task Selection을 얻지 못하면 기존 sched\_ext의
Local DSQ 선택 경로로 복귀한다.

또한 선택적으로 Fairness Guard를 활성화할 수 있다.
Fairness Guard는 Neural Policy가 선택한 Task와
Candidate 중 가장 작은 Weighted DSQ Virtual Time을 가지는
Task의 Virtual Time 차이를 비교한다.
이 차이가 설정된 Budget을 초과하면
Neural Policy의 선택 대신 해당 Fair Candidate를 실행한다.
이는 Neural Policy가 기존 CPU Service Ordering에서
과도하게 벗어나는 상황을 제한하기 위한 장치이며,
모든 Workload에서 Starvation이 발생하지 않음을 보장하는
절대적인 시간 제한은 아니다.

이와 같이 RL-SchedExt는 학습된 Policy를
기존 Scheduling Mechanism과 분리하여 적용하고,
유효한 State 및 Inference 결과가 확보된 경우에만
Neural Policy에 의한 Task Ordering을 수행한다.
